\section{A Nonlinear Caputo Double-Allee Ecological Realization}\label{sec:ecology}

We now ask whether the open matrix band in Equation~\eqref{eq:fractional-band} can be realized by a biologically standard nonlinear Kolmogorov system rather than by an abstract matrix construction. The answer depends on the interaction architecture. The first natural model fails for a structural reason; an intraguild-predation architecture removes precisely that obstruction.

\subsection{Positive-Equilibrium Orbit Equivalence}

Consider a commensurate Caputo Kolmogorov system
\begin{equation}
 {}^C D_t^\alpha x_i=x_iF_i(x),\qquad i=1,\dots,n,
 \label{eq:kolmogorov-general}
\end{equation}
and let \(x^*\gg0\) be an equilibrium. Put \(B=DF(x^*)\).

\begin{Lemma}[Kolmogorov orbit equivalence]\label{lem:kolmogorov-orbit}
At a positive equilibrium,
\begin{equation}
 J(x^*)=\diag(x^*)B.
 \label{eq:kolmogorov-jacobian}
\end{equation}
Consequently,
\[
 J(x^*)\in\mathcal F_\alpha^{(n)}
 \iff
 B\in\mathcal F_\alpha^{(n)},
 \qquad
 J(x^*)\in\mathcal D_H^{(n)}
 \iff
 B\in\mathcal D_H^{(n)}.
\]
\end{Lemma}

\begin{proof}
For \(G_i(x)=x_iF_i(x)\),
\[
 \partial_jG_i
 =
 \delta_{ij}F_i+x_i\partial_jF_i.
\]
At \(x^*\), \(F_i(x^*)=0\), which gives Equation~\eqref{eq:kolmogorov-jacobian}. Since \(D\mapsto D\diag(x^*)\) is a bijection of the positive diagonal cone onto itself, the two positive-diagonal spectral orbits coincide.
\end{proof}

This factorization is standard in Kolmogorov-system stability theory and is used here only as a bridge \cite{HouBaigent2015,ChenWangLiu2024}.

\subsection{A Structural No-Go for Two Competing Consumers}

Begin with a double-Allee basal prey consumed by two competitors:
\begin{align}
 {}^CD_t^\alpha x&=x[g_{\rm DA}(x)-q_1y-q_2z],\nonumber\\
 {}^CD_t^\alpha y&=y[e_1q_1x-\mu_1-c_{11}y-c_{12}z],\label{eq:competitive-model}\\
 {}^CD_t^\alpha z&=z[e_2q_2x-\mu_2-c_{21}y-c_{22}z],\nonumber
\end{align}
with positive parameters. At a positive equilibrium let \(s=-g_{\rm DA}'(X)>0\). The reduced matrix is
\begin{equation}
 B_{\rm comp}
 =
 \begin{pmatrix}
 -s&-q_1&-q_2\\
 e_1q_1&-c_{11}&-c_{12}\\
 e_2q_2&-c_{21}&-c_{22}
 \end{pmatrix}.
 \label{eq:Bcomp}
\end{equation}
On its strict-\(P\) stratum,
\begin{align}
 \beta_{12}&=1+\frac{e_1q_1^2}{sc_{11}},&
 \beta_{13}&=1+\frac{e_2q_2^2}{sc_{22}},&
 \beta_{23}&=1-\frac{c_{12}c_{21}}{c_{11}c_{22}}>0,
 \label{eq:competitive-betas}
\end{align}
and the normalized directed three-cycle contribution is
\begin{equation}
 L_3
 =
 \frac{q_1q_2(c_{12}e_2+c_{21}e_1)}
 {sc_{11}c_{22}}>0.
 \label{eq:competitive-L3}
\end{equation}
The determinant identity
\begin{equation}
 \kappa
 =
 \beta_{12}+\beta_{13}+\beta_{23}-2-L_3
 \label{eq:kappa-loop}
\end{equation}
then yields the following obstruction.

\begin{Theorem}[competitive-architecture no-go]\label{thm:no-go}
Every strict-\(P\) point of the model in Equation~\eqref{eq:competitive-model} is classically Hurwitz \(D\)-stable. In particular, the architecture cannot realize an open set satisfying
\[
 T_1(\beta)<\kappa<T_\alpha(\beta).
\]
\end{Theorem}

\begin{proof}
From Equations~\eqref{eq:competitive-L3}--\eqref{eq:kappa-loop},
\[
 \kappa
 <
 \beta_{12}+\beta_{13}+\beta_{23}-2
 <
 \beta_{12}+\beta_{13}+\beta_{23}.
\]
Because all \(\beta_{ij}>0\),
\[
 T_1(\beta)
 =
 \left(
 \sqrt{\beta_{12}}+\sqrt{\beta_{13}}+\sqrt{\beta_{23}}
 \right)^2
 >
 \beta_{12}+\beta_{13}+\beta_{23}.
\]
Thus \(\kappa<T_1(\beta)\), which is exactly Cain's strict-\(P\) \(D\)-stability inequality.
\end{proof}

The failure is therefore architectural rather than parametric: adding more parameter search cannot reverse the sign in Equation~\eqref{eq:competitive-L3}.

\subsection{Intraguild Predation and Exact Invariant Access}

We replace direct competition between the consumers by intraguild predation/omnivory, a standard three-species motif \cite{HoltPolis1997,BaiKangRuanWang2021,Panja2019}:
\begin{align}
 {}^CD_t^\alpha x
 &=
 x[g_{\rm DA}(x)-q_1y-q_2z],
 \nonumber\\
 {}^CD_t^\alpha y
 &=
 y[e_1q_1x-\mu_1-c_1y-hz],
 \label{eq:igp-model}\\
 {}^CD_t^\alpha z
 &=
 z[e_2q_2x+e_3hy-\mu_2-c_2z],
 \nonumber
\end{align}
where
\begin{equation}
 g_{\rm DA}(x)
 =
 \frac{r}{x+a}
 \left(1-\frac{x}{K}\right)(x-m),
 \qquad
 0<m<K,
 \label{eq:gDA}
\end{equation}
and all displayed parameters are positive. Here \(x\) is the basal prey/resource, \(y\) the intermediate consumer or intraguild prey, and \(z\) the omnivorous top predator.

At a positive equilibrium \((X,Y,Z)\),
\begin{align}
 g_{\rm DA}(X)&=q_1Y+q_2Z,\label{eq:eq1}\\
 c_1Y+hZ&=e_1q_1X-\mu_1,\label{eq:eq2}\\
 -e_3hY+c_2Z&=e_2q_2X-\mu_2.\label{eq:eq3}
\end{align}
Put
\[
 \Delta=c_1c_2+e_3h^2,
 \quad
 A_1=e_1q_1X-\mu_1,
 \quad
 A_2=e_2q_2X-\mu_2.
\]
Then
\begin{equation}
 Y=\frac{c_2A_1-hA_2}{\Delta},
 \qquad
 Z=\frac{e_3hA_1+c_1A_2}{\Delta}.
 \label{eq:YZ}
\end{equation}
Moreover
\begin{equation}
 q_1Y+q_2Z=\chi X-\nu,
 \label{eq:affine-predation}
\end{equation}
where
\begin{equation}
 \chi=
 \frac{
 e_1c_2q_1^2+e_2c_1q_2^2
 +hq_1q_2(e_1e_3-e_2)
 }{\Delta},
 \label{eq:chi}
\end{equation}
and
\[
 \nu=
 \frac{
 \mu_1(q_1c_2+q_2e_3h)
 +\mu_2(q_2c_1-q_1h)
 }{\Delta}.
\]
Thus \(X\) satisfies the degree-at-most-two equation
\begin{equation}
 (r+K\chi)X^2
 -
 [r(K+m)-K\chi a+K\nu]X
 +
 rKm-K\nu a
 =0.
 \label{eq:Xquadratic}
\end{equation}
In the constructive regime used below \(e_1e_3\ge e_2\), so \(\chi>0\) and Equation~\eqref{eq:Xquadratic} is a genuine quadratic.

At coexistence set
\[
 s=-g_{\rm DA}'(X).
\]
The reduced per-capita matrix is
\begin{equation}
 B=
 \begin{pmatrix}
 -s&-q_1&-q_2\\
 e_1q_1&-c_1&-h\\
 e_2q_2&e_3h&-c_2
 \end{pmatrix}.
 \label{eq:Bigp}
\end{equation}
Its order-two principal minors are
\[
 m_{12}=sc_1+e_1q_1^2,\quad
 m_{13}=sc_2+e_2q_2^2,\quad
 m_{23}=c_1c_2+e_3h^2,
\]
and a direct determinant calculation gives
\begin{equation}
 -\det B
 =
 \Delta(s+\chi).
 \label{eq:q-delta}
\end{equation}
Hence the exact strict-\(P\) condition for this architecture is
\begin{equation}
 s>0,\qquad s+\chi>0.
 \label{eq:igp-strictP}
\end{equation}
The four C-10 invariants are
\begin{align}
 \beta_{12}&=1+\frac{e_1q_1^2}{sc_1},&
 \beta_{13}&=1+\frac{e_2q_2^2}{sc_2},&
 \beta_{23}&=1+\frac{e_3h^2}{c_1c_2},
 \label{eq:igp-betas}\\
 \kappa
 &=
 \beta_{12}+\beta_{13}+\beta_{23}-2
 +
 \frac{hq_1q_2(e_1e_3-e_2)}{sc_1c_2}.
 \label{eq:igp-kappa}
\end{align}
Equivalently,
\begin{equation}
 L_3
 =
 \frac{hq_1q_2(e_2-e_1e_3)}{sc_1c_2}.
 \label{eq:igp-L3}
\end{equation}
Thus \(e_1e_3>e_2\) makes \(L_3<0\), the sign mechanism missing from the competitive architecture.

Figure~\ref{fig:mechanism} displays the two motifs and the corresponding directed-cycle sign. The diagram is a theorem schematic; no numerical data enter the no-go or realization argument.

\begin{figure}[H]
\centering
\includegraphics[width=0.94\textwidth]{fig05_ecological_mechanism.pdf}
\caption{Ecological architectures and the signed three-cycle mechanism. \textbf{(a)} With two competing consumers, both oriented three-cycles contribute with the sign that makes \(L_3>0\), forcing \(\kappa<T_1\) on the strict-\(P\) stratum. \textbf{(b)} Intraguild predation reverses one cycle; when \(e_1e_3>e_2\), \(L_3<0\) and access to \(\kappa>T_1\) becomes possible.}
\label{fig:mechanism}
\end{figure}

\begin{Theorem}[four-invariant IGP realization]\label{thm:igp-realization}
Fix
\[
 \beta_{12},\beta_{13},\beta_{23}>1
\]
and
\begin{equation}
 \kappa>
 \beta_{12}+\beta_{13}+\beta_{23}-2.
 \label{eq:igp-target}
\end{equation}
Then positive IGP parameters can be chosen so that the matrix in Equation~\eqref{eq:Bigp} has exactly these four invariants. The construction may be made with
\[
 0<e_1,e_2,e_3<1,
 \qquad
 e_1e_3>e_2.
\]
\end{Theorem}

\begin{proof}
Let
\[
 A=\beta_{12}-1,\quad
 B_0=\beta_{13}-1,\quad
 C=\beta_{23}-1,
\]
and
\[
 R=\kappa-(\beta_{12}+\beta_{13}+\beta_{23}-2)>0.
\]
Put
\[
 \rho=\frac{R}{\sqrt{AB_0C}},
 \qquad
 \tau=\frac{\rho+\sqrt{\rho^2+4}}{2}>1,
\]
so \(\tau-\tau^{-1}=\rho\). Choose any \(\eta\in(0,1)\) and set
\[
 e_1=e_3=\eta,
 \qquad
 e_2=\frac{\eta^2}{\tau^2}.
\]
For arbitrary \(s,c_1,c_2>0\), define
\begin{equation}
 q_1=\sqrt{\frac{Asc_1}{e_1}},
 \qquad
 q_2=\sqrt{\frac{B_0sc_2}{e_2}},
 \qquad
 h=\sqrt{\frac{Cc_1c_2}{e_3}}.
 \label{eq:igp-right-inverse}
\end{equation}
Equations~\eqref{eq:igp-betas} give the prescribed \(\beta\)'s, while
\[
 \frac{hq_1q_2(e_1e_3-e_2)}{sc_1c_2}
 =
 \sqrt{AB_0C}\left(\tau-\frac1\tau\right)
 =R.
\]
Equation~\eqref{eq:igp-kappa} then gives the prescribed \(\kappa\). Since \(\tau>1\), \(e_1e_3>e_2\).
\end{proof}

\subsection{Embedding the Invariants in the Double-Allee Model}

The previous theorem fixes the desired local prey self-slope \(s>0\). We now realize it at a positive equilibrium of Equation~\eqref{eq:igp-model}.

Fix
\[
 a>0,\qquad 0<m<X<K,
\]
and define
\begin{equation}
 H_A
 =
 \frac{1}{K-X}
 -
 \frac{1}{X-m}
 +
 \frac{1}{X+a}.
 \label{eq:HA}
\end{equation}
Writing
\[
 C_A=\frac{m+a}{(X-m)(X+a)},
\]
we have \(H_A=(K-X)^{-1}-C_A\), so \(H_A>0\) precisely when
\begin{equation}
 X<K<
 K_{\max}
 :=
 X+\frac{(X-m)(X+a)}{m+a}.
 \label{eq:Kmax}
\end{equation}
Set
\[
 Q=\frac{s}{H_A}.
\]
Choose \(\xi\in(0,1)\) and prescribe
\begin{equation}
 Y=\frac{\xi Q}{q_1},
 \qquad
 Z=\frac{(1-\xi)Q}{q_2}.
 \label{eq:YZsplit}
\end{equation}
Then \(q_1Y+q_2Z=Q\). Define
\begin{align}
 \mu_1&=e_1q_1X-c_1Y-hZ,\label{eq:mu1}\\
 \mu_2&=e_2q_2X+e_3hY-c_2Z.\label{eq:mu2}
\end{align}
The first mortality is positive whenever
\[
 Q<Q_1:=
 \frac{e_1q_1X}
 {\xi c_1/q_1+(1-\xi)h/q_2}.
\]
If \((1-\xi)c_2/q_2-\xi e_3h/q_1>0\), the second is positive whenever
\[
 Q<Q_2:=
 \frac{e_2q_2X}
 {(1-\xi)c_2/q_2-\xi e_3h/q_1};
\]
otherwise \(\mu_2>0\) automatically. Let \(\bar Q=Q_1\) in the latter case and \(\bar Q=\min(Q_1,Q_2)\) in the former. It is enough to choose
\begin{equation}
 0<K-X<
 \frac{1}{C_A+s/\bar Q}.
 \label{eq:K-close-bound}
\end{equation}
Finally set
\begin{equation}
 r=
 \frac{Q(X+a)}
 {(1-X/K)(X-m)}.
 \label{eq:r-embedding}
\end{equation}
Then \(g_{\rm DA}(X)=Q=q_1Y+q_2Z\), while
\[
 g_{\rm DA}'(X)
 =
 Q\left[
 -\frac1{K-X}
 +\frac1{X-m}
 -\frac1{X+a}
 \right]
 =-QH_A=-s.
\]
Thus the positive equilibrium has exactly the reduced matrix constructed in Theorem~\ref{thm:igp-realization}.

\begin{Theorem}[open genuinely fractional biological region]\label{thm:open-ecological}
For every fixed \(0<\alpha<1\), there exists a nonempty open set in the full positive biological parameter space of Equation~\eqref{eq:igp-model} for which a positive coexistence equilibrium on the constructed smooth branch satisfies
\begin{equation}
 J\in
 \mathcal F_\alpha^{(3)}
 \setminus
 \mathcal D_H^{(3)}.
 \label{eq:open-ecological-band}
\end{equation}
\end{Theorem}

\begin{proof}
If \(2/3<\alpha<1\), choose \(\beta\in(1,\infty)^3\) and
\[
 T_1(\beta)<\kappa<T_\alpha(\beta).
\]
Theorem~\ref{thm:igp-realization} gives an IGP reduced matrix with these invariants, and Equations~\eqref{eq:HA}--\eqref{eq:r-embedding} embed it at a positive double-Allee coexistence equilibrium. Theorem~\ref{thm:c10} gives \(B\in\mathcal F_\alpha^{(3)}\setminus\mathcal D_H^{(3)}\), and Lemma~\ref{lem:kolmogorov-orbit} transfers the statement to the actual Jacobian \(J\).

If \(0<\alpha\le2/3\), choose \(\beta>1\) and any \(\kappa>T_1(\beta)\) satisfying Equation~\eqref{eq:igp-target}. The construction is strict-\(P\), so Kellogg's wedge implies \(B\in\mathcal F_\alpha^{(3)}\), while Cain's inequality fails.

All feasibility and stability inequalities are strict. Moreover
\[
 -\det B=\kappa sc_1c_2>0,
\]
so the state Jacobian of the coexistence equations is nonsingular. The implicit-function theorem preserves the positive equilibrium on a smooth local branch under sufficiently small perturbations of all biological parameters, and continuity preserves the strict invariant inequalities. Hence the set is open in the full biological parameter space.
\end{proof}

\subsection{The Allee Threshold as a Monotone Invariant-Space Control}

After eliminating \(Y\) and \(Z\), the coexistence equation is
\begin{equation}
 F(X,m)
 :=
 g_{\rm DA}(X;m)-\chi X+\nu=0.
 \label{eq:scalar-equilibrium}
\end{equation}
Along a positive strict-\(P\) branch put
\[
 Q=g_{\rm DA}(X)=\chi X-\nu>0,
 \qquad
 s=-g_{\rm DA}'(X)>0.
\]
By Equation~\eqref{eq:q-delta},
\[
 s+\chi>0.
\]
Implicit differentiation gives the first exact monotonicity law:
\begin{equation}
 \frac{dX}{dm}
 =
 -\frac{Q}{(X-m)(s+\chi)}
 <0.
 \label{eq:dXdm}
\end{equation}

The slope \(s\) decreases as well. Set
\[
 d=X-m,\qquad
 \ell=K-X,\qquad
 w=X+a.
\]
Then \(H_A=\ell^{-1}-d^{-1}+w^{-1}\), and direct algebra gives
\begin{equation}
 \partial_XH_A-H_A^2
 =
 \frac{2(w-d)(\ell+w)}{d\ell w^2}>0,
 \label{eq:H-identity}
\end{equation}
because \(w-d=m+a>0\). Differentiating \(s=QH_A\) along Equation~\eqref{eq:scalar-equilibrium} yields
\begin{equation}
 \frac{ds}{dm}
 =
 -\frac{Q}{d(s+\chi)}
 \left[
 Q\partial_XH_A+\frac{s}{d}
 +\chi\left(H_A+\frac1d\right)
 \right].
 \label{eq:dsdm}
\end{equation}
Since strict-\(P\) gives \(\chi>-s\), the bracket is strictly larger than
\[
 Q\partial_XH_A+\frac{s}{d}
 -s\left(H_A+\frac1d\right)
 =
 Q(\partial_XH_A-H_A^2)>0.
\]
Therefore
\begin{equation}
 \boxed{
 \frac{dX}{dm}<0,\qquad
 \frac{ds}{dm}<0
 }
 \label{eq:allee-monotonicity}
\end{equation}
throughout the positive strict-\(P\) coexistence stratum. No sign assumption on \(\chi\) is required.

Put \(t=1/s\) and define
\[
 A_0=\frac{e_1q_1^2}{c_1},\quad
 B_0=\frac{e_2q_2^2}{c_2},\quad
 C_0=1+\frac{e_3h^2}{c_1c_2},\quad
 E_0=\frac{hq_1q_2(e_1e_3-e_2)}{c_1c_2}.
\]
Then
\begin{equation}
 \beta_{12}=1+A_0t,\quad
 \beta_{13}=1+B_0t,\quad
 \beta_{23}=C_0,\quad
 \kappa=C_0+(A_0+B_0+E_0)t,
 \label{eq:t-path}
\end{equation}
and Equation~\eqref{eq:allee-monotonicity} implies \(dt/dm>0\).

Along this path the classical gap is
\begin{align}
 G_1(t)
 &:=
 \kappa(t)-T_1(\beta(t))
 \nonumber\\
 &=
 E_0t-2
 -2\sqrt{(1+A_0t)(1+B_0t)}
 \nonumber\\
 &\quad
 -2\sqrt{C_0}
 \left[
 \sqrt{1+A_0t}
 +
 \sqrt{1+B_0t}
 \right].
 \label{eq:G1}
\end{align}

\begin{Theorem}[unique transverse Cain crossing]\label{thm:crossing}
If
\begin{equation}
 E_0>2\sqrt{A_0B_0},
 \label{eq:cross-condition}
\end{equation}
then \(G_1\) has a unique positive zero \(t_H\). The crossing is transverse and satisfies
\begin{equation}
 G_1'(t_H)
 \ge
 \frac{4+4\sqrt{C_0}}{t_H}
 >0.
 \label{eq:transversality}
\end{equation}
Consequently, because \(t\) increases strictly with \(m\), the coexistence branch crosses the classical \(D\)-stability boundary at most once and does so transversally.
\end{Theorem}

\begin{proof}
We have
\[
 G_1(0)=-4-4\sqrt{C_0}<0,
\]
whereas Equation~\eqref{eq:cross-condition} makes the asymptotic linear coefficient positive, so \(G_1(t)\to+\infty\). Direct differentiation shows
\begin{align}
 G_1''(t)
 &=
 \frac{(A_0-B_0)^2}
 {2[(1+A_0t)(1+B_0t)]^{3/2}}
 \nonumber\\
 &\quad+
 \sqrt{C_0}\left[
 \frac{A_0^2}{2(1+A_0t)^{3/2}}
 +
 \frac{B_0^2}{2(1+B_0t)^{3/2}}
 \right]>0.
\end{align}
Thus \(G_1\) is strictly convex and has exactly one positive root. Convexity gives
\[
 G_1'(t_H)
 \ge
 \frac{G_1(t_H)-G_1(0)}{t_H},
\]
which is Equation~\eqref{eq:transversality}.
\end{proof}

The crossing can be prescribed. For any \(t_0>0\), set
\begin{equation}
 E_0
 =
 \frac{
 2+2\sqrt{(1+A_0t_0)(1+B_0t_0)}
 +2\sqrt{C_0}\left[
 \sqrt{1+A_0t_0}+\sqrt{1+B_0t_0}
 \right]
 }{t_0}.
 \label{eq:prescribed-E0}
\end{equation}
Then \(G_1(t_0)=0\) and Equation~\eqref{eq:cross-condition} holds. Theorem~\ref{thm:igp-realization} realizes this \(E_0\) through a joint parameter choice, and the double-Allee embedding independently sets \(s_0=1/t_0\) at a selected \(m_0\). Hence, with every non-\(m\) model parameter fixed, there exists \(\varepsilon>0\) such that
\[
 m_0-\varepsilon<m<m_0
 \quad\Longrightarrow\quad
 J\in\mathcal D_H^{(3)},
\]
whereas
\begin{equation}
 m_0<m<m_0+\varepsilon
 \quad\Longrightarrow\quad
 J\in
 \mathcal F_\alpha^{(3)}
 \setminus
 \mathcal D_H^{(3)}.
 \label{eq:m-transition}
\end{equation}
For \(\alpha>2/3\), the right-hand implication follows from \(T_\alpha>T_1\) and continuity; for \(\alpha\le2/3\), strict-\(P\) itself gives fractional \(D\)-stability.

Figure~\ref{fig:m-path} visualizes this mechanism for the certified realization used in Section~\ref{sec:certified}. All non-\(m\) biological parameters are held fixed.

\begin{figure}[H]
\centering
\includegraphics[width=0.95\textwidth]{fig06_double_allee_m_path.pdf}
\caption{The Allee threshold drives a fixed three-species community through exact invariant space at \(\alpha=0.9\). \textbf{(a)} Cain's gap \(\kappa-T_1\) changes sign once, at the exact crossing \(m_0=0.374451\ldots\), while \(T_{0.9}-\kappa\) remains positive on the displayed fractional-only branch. \textbf{(b)} The same path in an invariant slice. The curves use exact formulas; the marked anchor is interval-certified in Section~\ref{sec:certified}.}
\label{fig:m-path}
\end{figure}

\subsection{Exact Two- versus Three-Dimensional Ecological Contrast}

For comparison, consider the two-species model
\[
 {}^CD_t^\alpha x=x[g_{\rm DA}(x)-qy],
 \qquad
 {}^CD_t^\alpha y=y[px-d].
\]
At coexistence \(X=d/p\), the reduced matrix is
\[
 B_2=
 \begin{pmatrix}
 g_{\rm DA}'(X)&-q\\
 p&0
 \end{pmatrix}.
\]
Theorem~\ref{thm:2d} gives
\[
 B_2\in\mathcal F_\alpha^{(2)}
 \iff
 g_{\rm DA}'(X)\le0,
\]
whereas classical Hurwitz \(D\)-stability requires \(g_{\rm DA}'(X)<0\). Hence the genuinely fractional difference occurs only on the codimension-one condition
\[
 g_{\rm DA}'(X)=0.
\]
Solving this condition for the Allee threshold gives
\begin{equation}
 m_c
 =
 \frac{X^2+2aX-Ka}{K+a}.
 \label{eq:mc-2d}
\end{equation}
Thus the two-species mechanism is exceptional, while Theorem~\ref{thm:open-ecological} gives a nonempty open fractional-only region in the three-species architecture.
